\section*{Hoofdstuk \ref{ch:b1:weathering-clays} --- Verwering en kleien}

\begin{solution}{exo:b1:weathering-clays:1}
\ce{K+}: $1/1.38 = 0.72$; \ce{Fe^2+}: $2/0.78 = 2.56$; beide onder 3, oplosbare
kationen. \ce{Fe^3+}: $3/0.645 = 4.65$; \ce{Ti^4+}: $4/0.605 = 6.61$; beide tussen
3 en 10, hydrolysaten. In een waterdoorzadigde bodem, zonder zuurstof, wordt ijzer
vrijgemaakt als \ce{Fe^2+}, dat in oplossing blijft en door het water wordt
meegevoerd; in een goed gedraineerde bodem maakt zuurstof er \ce{Fe^3+} van, dat meteen
als oxiden en hydroxiden neerslaat en blijft.
\end{solution}

\begin{solution}{exo:b1:weathering-clays:2}
A zwelt met glycol en klapt in bij verhitting: een \omterm{def:b1:weathering-clays:clay-mineral}{smectiet}. B ligt op
\SI{7}{\angstrom} en wordt vernield bij \SI{550}{\celsius}: \omterm{def:b1:weathering-clays:clay-mineral}{kaoliniet}. C, op
\SI{10}{\angstrom} wat de behandeling ook is: \omterm{def:b1:weathering-clays:clay-mineral}{illiet}. D, op \SI{14}{\angstrom} zonder
zwelling of inklappen: \omterm{def:b1:weathering-clays:clay-mineral}{chloriet}.
\end{solution}

\begin{solution}{exo:b1:weathering-clays:3}
Twee orthoklazen leveren het 2 aluminium van één \omterm{def:b1:weathering-clays:clay-mineral}{kaoliniet}; $6 - 2 = 4$ silicium
verlaat als \ce{H4SiO4}; 2 \ce{K+} verlaten, in balans met 2 \ce{H+}; zuurstof,
$16 + w = 9 + 16$, geeft $w = 9$:
\[
  \ce{2 KAlSi3O8 + 2 H+ + 9 H2O -> Al2Si2O5(OH)4 + 2 K+ + 4 H4SiO4}.
\]
Waterstof klopt: $2 + 18 = 20 = 4 + 16$. Al het kalium verlaat, en $4/6 =
2/3$ van het silicium.
\end{solution}

\begin{solution}{exo:b1:weathering-clays:4}
(a) De massa verandert niet: $10 \times 999.87/916.72 = \qty{10.91}{cm^3}$.
(b) Tussen ongeveer \qty{-4}{\celsius} en \qty{-15}{\celsius}. Net onder nul kan het ijs
niet genoeg druk opbouwen om de scheur open te wiggen; bij strengere vorst beweegt het
water niet meer door de poriën naar het ijs toe, dat ophoudt te groeien. (c) Volgens
\cref{prop:b1:weathering-clays:surface} met $n = \qty{1}{m}/\qty{1}{cm} = 100$
wordt het oppervlak met 100 vermenigvuldigd: van \qty{6}{m^2} tot \qty{600}{m^2}.
\end{solution}

\begin{solution}{exo:b1:weathering-clays:5}
$S/m = 6/(\rho d)$: voor \qty{1}{cm}, $6/(2600\times0.01) =
\qty{0.231}{m^2.kg^{-1}}$; voor \qty{1}{mm}, \qty{2.31}{m^2.kg^{-1}}; voor
\qty{1}{\micro m}, \qty{2308}{m^2.kg^{-1}}, of \qty{2.3}{m^2.g^{-1}}.
Chemische reacties vinden plaats aan het oppervlak van de korrels, dus groeit hun
snelheid per kilogram gesteente met het specifieke oppervlak, tienduizend keer van de
centimeter tot de micrometer. \omterm{def:b1:weathering-clays:weathering}{Fysische verwering} snijdt de rots in steeds kleinere
stukken en opent scheuren waardoor water binnendringt: zij vermenigvuldigt het oppervlak
waarop \omterm{def:b1:weathering-clays:weathering}{chemische verwering} kan werken.
\end{solution}

\begin{solution}{exo:b1:weathering-clays:6}
Een ton albiet is $\num{e6}/262.22 = \qty{3813.6}{mol}$. Volgens
\cref{prop:b1:weathering-clays:feldspars} geven twee mol één mol \omterm{def:b1:weathering-clays:clay-mineral}{kaoliniet}:
$\qty{1906.8}{mol} \times \qty{258.16}{g.mol^{-1}} = \qty{492.3}{kg}$. Twee van de drie
silicium-atomen verlaten: $2\times3813.6 = \qty{7627.2}{mol}$ van
\ce{SiO2}, \qty{458.2}{kg}. Al het natrium verlaat: $3813.6\times22.99 =
\qty{87.7}{kg}$. Minder dan de helft van de massa van het \omterm{def:b1:rock-forming-minerals:groups}{veldspaat} blijft als klei over.
\end{solution}

\begin{solution}{exo:b1:weathering-clays:7}
Volgens \cref{prop:b1:weathering-clays:immobile} geeft een kilogram granit
$0.30/0.60 = \qty{0.5}{kg}$ saproliet. Verloren deel: natrium $1 -
(0.35/3.5)(0.5) = 0.95$; silicium $1 - (60/72)(0.5) = 0.583$; aluminium $1 -
(27/14)(0.5) = 0.036$. Aluminium blijft vrijwel behouden, het natrium is bijna volledig
weg, en bijna \qty{60}{\percent} van het silicium ook, hoewel veel van wat overblijft in
kwarts zit: de veldspaten worden in \omterm{def:b1:weathering-clays:clay-mineral}{kaoliniet} omgezet, \omterm{def:b1:weathering-clays:hydrolysis-degree}{monosiallitisatie}.
\end{solution}

\begin{solution}{exo:b1:weathering-clays:8}
Pyrofylliet: $x = y = 0$, lading 0. Muscovietlaag: $x = 1$, lading $-1$.
Beidelliet-laag: $x = 0.5$, lading $-0.5$. Montmorilloniet-laag: $y = 0.33$,
lading $-0.33$. De twee smectietlagen, beidelliet en montmorilloniet, kunnen
zwellen: hun kleine lading wordt door weinig tussenlaagkationen in balans gehouden, die
de lagen los houden en water aantrekken terwijl zij gehydrateerd raken. Pyrofylliet heeft
geen tussenlaagkationen om te hydrateren, en de lading $-1$ van muscoviet wordt in balans
gehouden door kaliumionen die de lagen vastzetten.
\end{solution}

\begin{solution}{exo:b1:weathering-clays:9}
(a) Mg: 2 en 2; Si: 1 en 1; C: 4 en 4; O: $4 + 8 + 4 = 16$ en $12 + 4 =
16$; H: 8 en $4 + 4 = 8$; lading: 0 en $4 - 4 = 0$. (b) De molaire massa van
forsteriet is $2\times24.305 + 28.085 + 4\times15.999 =
\qty{140.69}{g.mol^{-1}}$, dus een ton is \qty{7108}{mol}. Zij verbruikt
$4\times7108 = \qty{28431}{mol}$ van \ce{CO2}, \qty{1251}{kg}; met twee
magnesium per formule-eenheid slaat \cref{prop:b1:weathering-clays:co2}
$2\times7108$ mol op, \qty{626}{kg}, de helft van het verbruikte. (c) Een ton calciet,
\qty{9991}{mol}, verbruikt \qty{9991}{mol} van \ce{CO2} (\qty{440}{kg}) bij het oplossen,
en geeft het allemaal terug wanneer het carbonaat opnieuw neerslaat: er wordt niets
opgeslagen.
\end{solution}

\begin{solution}{exo:b1:weathering-clays:10}
(a) $3 < 3/r < 10$ geeft $\qty{0.3}{\angstrom} < r < \qty{1.0}{\angstrom}$.
(b) $3/0.27 = 11.1 > 10$: boor reist als oxyanion, boraat. (c) $2/0.83 =
2.41$, een oplosbaar kation; $4/0.53 = 7.55$, een hydrolysaat. In water arm aan
zuurstof is mangaan opgelost als \ce{Mn^2+}; waar de bron de lucht ontmoet, wordt het
geoxideerd tot \ce{Mn^4+}, dat meteen als een zwarte oxide op de rots neerslaat. (d) De
\omterm{def:b1:weathering-clays:ionic-potential}{ionische potentiaal}, $4/0.40 = 10.0$, ligt op de grens tussen hydrolysaten en oxyanionen:
silicium wordt niet vastgehouden zoals aluminium en ook niet meegevoerd zoals zwavel, en
het uitgeloogde deel hangt af van hoeveel water er doorheen gaat, de basis van de drie
graden van \omterm{def:b1:weathering-clays:hydrolysis}{hydrolyse}.
\end{solution}

\begin{solution}{exo:b1:weathering-clays:11}
(a) Met \qty{12.011}{g.mol^{-1}} koolstof is \num{11.7e12} tot
\qty{13.6e12}{mol} 0.141 tot 0.163 miljard ton koolstof per jaar; met
\qty{44.009}{g.mol^{-1}}, 0.515 tot 0.599 miljard ton koolstofdioxide.
(b) $4.08/13.6 = \qty{30}{\percent}$ tot $4.08/11.7 = \qty{35}{\percent}$.
(c) De helft van het verbruikte koolstofdioxide wordt opgeslagen wanneer het calcium en
magnesium als carbonaat neerslaan: 0.070 tot 0.082 miljard ton koolstof per jaar. Het
werkelijke cijfer is kleiner, omdat een deel van het bicarbonaat met natrium en kalium
meegaat, die op de zeebodem geen carbonaat vormen.
\end{solution}

\begin{solution}{exo:b1:weathering-clays:12}
Si/Al $= (w_{\ce{SiO2}}/60.08)/(2w_{\ce{Al2O3}}/101.96)$. A: $0.7748/0.7748 =
1.00$, de verhouding van \omterm{def:b1:weathering-clays:clay-mineral}{kaoliniet}, \omterm{def:b1:weathering-clays:hydrolysis-degree}{monosiallitisatie}. B: $0.9987/0.4904 = 2.04$, een
2:1-kleimineraal, \omterm{def:b1:weathering-clays:hydrolysis-degree}{bisiallitisatie}. C: $0.0333/1.2358 = 0.027$, bijna geen silicium meer
over, gibbsiet, \omterm{def:b1:weathering-clays:hydrolysis-degree}{allitisatie}. Van de minst naar de intensiefste uitloging, dat is naar het
heetste en natste klimaat toe: B, A, C.
\end{solution}

\begin{solution}{pb:b1:weathering-clays:1}
\textbf{1.} \ce{2 NaAlSi3O8 + 2 H+ + 9 H2O -> Al2Si2O5(OH)4 + 2 Na+ + 4 H4SiO4}.

\textbf{2.} \ce{CaAl2Si2O8 + 2 H+ + H2O -> Al2Si2O5(OH)4 + Ca^2+}. Geen silicium
verlaat: anorthiet heeft al het ene silicium per aluminium van \omterm{def:b1:weathering-clays:clay-mineral}{kaoliniet}.

\textbf{3.} Eén gibbsiet neemt één aluminium op, dus een half anorthiet: verdubbel de reactie,
\ce{CaAl2Si2O8 + 2 H+ + 6 H2O -> 2 Al(OH)3 + Ca^2+ + 2 H4SiO4}. Zuurstof:
$8 + 6 = 6 + 8$; waterstof: $2 + 12 = 6 + 8$.

\textbf{4.} \ce{5 NaAlSi3O8 + 4 H+ + 16 H2O -> 2 Na_{0.5}Al2(Si_{3.5}Al_{0.5})O10(OH)2 + 4 Na+ + 8 H4SiO4}; % ce-unbalanced-ok: decimal subscripts the checker cannot read; balance tested in tests/bachelor-1/test_weathering_clays.py
waterstof, $4 + 32 = 36$ aan de linkerkant en $2\times2 + 8\times4 = 36$ aan de
rechterkant, klopt.

\textbf{5.} Elke \ce{H+} gaat samen met één \ce{HCO3-} uit één \ce{CO2} en één
\ce{H2O}:
\ce{2 NaAlSi3O8 + 2 CO2 + 11 H2O -> Al2Si2O5(OH)4 + 2 Na+ + 2 HCO3- + 4 H4SiO4}
\ce{CaAl2Si2O8 + 2 CO2 + 3 H2O -> Al2Si2O5(OH)4 + Ca^2+ + 2 HCO3-}.

\textbf{6.} Albiet naar beidelliet is \omterm{def:b1:weathering-clays:hydrolysis-degree}{bisiallitisatie}; albiet en anorthiet naar
\omterm{def:b1:weathering-clays:clay-mineral}{kaoliniet} zijn \omterm{def:b1:weathering-clays:hydrolysis-degree}{monosiallitisatie}; anorthiet naar gibbsiet is \omterm{def:b1:weathering-clays:hydrolysis-degree}{allitisatie}.

\textbf{7.} $\num{e6}/262.22 = \qty{3813.6}{mol}$ albiet;
$\num{e6}/278.20 = \qty{3594.5}{mol}$ anorthiet.

\textbf{8.} Albiet: $3813.6/2 = \qty{1906.8}{mol}$ \omterm{def:b1:weathering-clays:clay-mineral}{kaoliniet}, \qty{492.3}{kg}.
Anorthiet: \qty{3594.5}{mol}, \qty{928.0}{kg}.

\textbf{9.} Albiet: $2\times3813.6 = \qty{7627.2}{mol}$ \ce{SiO2}, \qty{458.2}{kg}, en
\qty{3813.6}{mol} natrium, \qty{87.7}{kg}. Anorthiet: geen silica, en
\qty{3594.5}{mol} calcium, \qty{144.1}{kg}.

\textbf{10.} Albiet: $1/2.62 = \qty{0.382}{m^3}$; zijn \omterm{def:b1:weathering-clays:clay-mineral}{kaoliniet},
$0.4923/2.6 = \qty{0.189}{m^3}$: het vaste stof verliest \qty{50.4}{\percent} van zijn
volume. Anorthiet: $1/2.75 = \qty{0.364}{m^3}$; zijn \omterm{def:b1:weathering-clays:clay-mineral}{kaoliniet},
$0.9280/2.6 = \qty{0.357}{m^3}$: een verlies van slechts \qty{1.8}{\percent}.

\textbf{11.} De achtergebleven ruimte wordt poriën, gevuld met water en lucht: een
saproliet op een albietrijk granit is licht, poreus en broos, met zijn kwartskorrels los
in een beetje klei. Anorthiet wordt bijna zijn eigen volume aan \omterm{def:b1:weathering-clays:clay-mineral}{kaoliniet}, dus de
verweerde gabbro behoudt het meeste van zijn vaste stof.

\textbf{12.} Natrium ($z/r = 0.98$) en calcium (2.0) zijn oplosbare kationen en verlaten;
aluminium (5.6) is een hydrolysaat en blijft geheel in de \omterm{def:b1:weathering-clays:clay-mineral}{kaoliniet}; silicium (10.0), op
de grens, verlaat gedeeltelijk uit albiet en helemaal niet uit anorthiet, waarvan het
silicium \omterm{def:b1:weathering-clays:clay-mineral}{kaoliniet} kan vasthouden.

\textbf{13.} Één \ce{CO2} per albiet: \qty{3813.6}{mol}, $3813.6 \times 44.01 =
\qty{167.8}{kg}$.

\textbf{14.} Twee per anorthiet: \qty{7189.0}{mol}, \qty{316.4}{kg}.

\textbf{15.} $7189.0 \times 61.02 = \qty{438.7}{kg}$ bicarbonaat.

\textbf{16.} Organismen slaan het neer als calciumcarbonaat,
\ce{Ca^2+ + 2 HCO3- -> CaCO3 + CO2 + H2O}: voor elk calcium gaat één van de twee
\ce{CO2} terug naar de lucht, \qty{3594.5}{mol} of \qty{158.2}{kg}, en de andere wordt
in het carbonaat begraven.

\textbf{17.} Natrium vormt op de zeebodem geen carbonaat: zijn bicarbonaat kan niet als
carbonaatgesteente worden opgeslagen, en de opslag op lange termijn wordt alleen uit
calcium en magnesium geteld.

\textbf{18.} Een ton calciet verbruikt \qty{9991}{mol} \ce{CO2} wanneer het oplost, en
geeft het allemaal terug wanneer het neerslaat: er wordt niets opgeslagen, en alleen de
verwering van calcium- en magnesiumsilicaten trekt koolstofdioxide voorgoed omlaag.

\textbf{19.} $\qty{200}{m^3.s^{-1}} \times \qty{0.3}{mol.m^{-3}} =
\qty{60}{mol.s^{-1}}$; over een jaar van \qty{3.156e7}{s}, \qty{1.893e9}{mol}, of
$\num{1.893e9}/\num{50000} = \qty{37870}{mol}$ per vierkante kilometer.

\textbf{20.} Twee bicarbonaten per anorthiet: \qty{9.467e8}{mol} per jaar,
$\num{9.467e8}\times\qty{278.20}{g} = \qty{263400}{t}$ per jaar, of
\qty{5.27}{t} per vierkante kilometer.

\textbf{21.} Het verse gesteente dat wordt verbruikt is $5.27/0.20 = \qty{26.3}{t}$, of
$\num{26340}/2700 = \qty{9.75}{m^3}$, per vierkante kilometer en per jaar: een laag van
\qty{9.75}{\micro m} dik per jaar, \qty{9.8}{m} per miljoen jaar.

\textbf{22.} Eén \ce{CO2} opgeslagen per calcium: $\num{9.467e8} \times 44.01 =
\qty{4.17e10}{g}$, ongeveer \qty{41700}{t} per jaar.

\textbf{23.} Het telt alleen de materie die in oplossing wordt meegevoerd. De \omterm{def:b1:weathering-clays:clay-mineral}{kaoliniet} en
het kwarts die door verwering achterblijven, worden ook door erosie verwijderd, als modder
en zand in de rivier, en worden niet geteld; evenmin is de verwering van de andere
mineralen van de rots meegerekend. En als een deel van het bicarbonaat uit natrium- of
kaliumveldspaten komt, die één bicarbonaat per formule-eenheid geven in plaats van twee,
moest meer gesteente verweren om het te leveren.

\textbf{24.} Volgens \cref{prop:b1:weathering-clays:co2} wordt één mol \ce{CO2} per mol
anorthiet opgeslagen: $3594.5 \times \qty{44.01}{g} = \qty{158.19}{kg}$. De verwering van
één ton anorthiet onttrekt voorgoed \textbf{158}~kg koolstofdioxide.
\end{solution}
